\documentclass[11pt,a4paper]{article}
\usepackage[utf8]{inputenc}
\usepackage[T1]{fontenc}
\usepackage[catalan]{babel}
\usepackage{lmodern}
\usepackage[sfdefault,lining]{FiraSans}
\usepackage[a4paper,top=2.3cm,bottom=2.2cm,left=1.8cm,right=1.8cm,headheight=1.3cm,headsep=0.4cm,footskip=1cm]{geometry}
\usepackage{amsmath,amssymb}
\usepackage{icomma}
\usepackage{graphicx}
\usepackage[table]{xcolor}
\usepackage{tikz}
\usepackage{fontawesome5}
\usepackage{ccicons}
\usepackage{fancyhdr}
\usepackage{lastpage}
\usepackage{tasks}
\usepackage{needspace}
\usepackage[most]{tcolorbox}
\usepackage{array}
\usepackage[hidelinks,pdftitle={Fitxa d'exercicis de la unitat 6: Proporcionalitat i percentatges. 2n d'ESO},pdfsubject={Matemàtiques 2n ESO},pdfkeywords={raons, proporcions, proporcionalitat directa, proporcionalitat inversa, regla de tres, proporcionalitat composta, repartiments proporcionals, percentatges, augments i disminucions percentuals}]{hyperref}

\definecolor{navy}{HTML}{1B2748}
\definecolor{accent}{HTML}{E8603A}
\definecolor{teal}{HTML}{1F9E90}
\definecolor{soft}{HTML}{EEF1F8}
\definecolor{muted}{HTML}{5D6685}

\setlength{\parindent}{0pt}
\setlength{\parskip}{3pt}
\settasks{label=\alph*), label-format=\bfseries\color{navy}, label-width=1.3em, item-indent=1.8em, column-sep=0.6em, after-item-skip=5pt, before-skip=3pt, after-skip=4pt}

% ---------- capçalera i peu ----------
\newcommand{\llicencia}{\href{https://creativecommons.org/licenses/by-nc-sa/4.0/deed.ca}{\ccbyncsa}\ \ Llicència Creative Commons Reconeixement-NoComercial-CompartirIgual 4.0 Internacional (CC BY-NC-SA 4.0)}
\pagestyle{fancy}
\fancyhf{}
\fancyhead[L]{\raisebox{-0.35\height}{\includegraphics[height=0.95cm]{logo-gva.png}}}
\fancyhead[C]{\small\color{muted}Fitxa d'exercicis · Unitat 6: Proporcionalitat i percentatges · 2n d'ESO}
\fancyhead[R]{\raisebox{-0.35\height}{\includegraphics[height=0.95cm]{logo-vermellar.png}}}
\fancyfoot[L]{\scriptsize\color{muted}\llicencia}
\fancyfoot[R]{\small\color{muted}\thepage}
\renewcommand{\headrulewidth}{0.4pt}
\renewcommand{\headrule}{\color{muted!40}\hrule height \headrulewidth}
\fancypagestyle{primera}{\fancyhf{}\renewcommand{\headrulewidth}{0pt}
  \fancyfoot[L]{\scriptsize\color{muted}\llicencia}
  \fancyfoot[R]{\small\color{muted}\thepage}}

% ---------- nivells i exercicis ----------
\newcommand{\estrelles}[1]{{\color{accent}\ifcase#1\or\faStar\faStar[regular]\faStar[regular]\or\faStar\faStar\faStar[regular]\or\faStar\faStar\faStar\fi}}
\newcounter{exer}
\newcommand{\ex}[1]{\par\needspace{4\baselineskip}\medskip\stepcounter{exer}\noindent{\color{navy}\bfseries\theexer.}\ \ {\small\estrelles{#1}}\ \ }

\newcommand{\seccio}[2]{\par\needspace{10\baselineskip}\bigskip
  \begin{tcolorbox}[colback=navy,colframe=navy,coltext=white,arc=2mm,boxrule=0pt,left=3mm,right=3mm,top=1.2mm,bottom=1.2mm]
  {\large\bfseries #1\quad #2}
  \end{tcolorbox}\vspace{-2pt}}

\newcolumntype{C}[1]{>{\centering\arraybackslash}p{#1}}
\usepackage{pifont}
% quadre de recordatori (només a la fitxa d'exercicis)
\newcommand{\recorda}[1]{\par\smallskip\begin{tcolorbox}[colback=accent!7,colframe=accent!55,arc=2mm,boxrule=0.5pt,left=2.5mm,right=2.5mm,top=1mm,bottom=1mm]\small{\color{accent}\faLightbulb[regular]}\ \textbf{Recorda:}\ #1\end{tcolorbox}}
\newcommand{\res}[1]{#1}\newcommand{\rest}[1]{#1}

\def\fmneed{14}
% =====================================================================
% Capçalera FentMates de les fitxes: un color per a cada curs.
% Els colors són els de la pàgina inicial (LEVELS a assets/js/materials.js);
% el logo ve de _eines/logo/pdf/ (python3 _eines/logo/genera.py).
%
% Ús, al final del preàmbul de la unitat:
%   % =====================================================================
% Capçalera FentMates de les fitxes: un color per a cada curs.
% Els colors són els de la pàgina inicial (LEVELS a assets/js/materials.js);
% el logo ve de _eines/logo/pdf/ (python3 _eines/logo/genera.py).
%
% Ús, al final del preàmbul de la unitat:
%   % =====================================================================
% Capçalera FentMates de les fitxes: un color per a cada curs.
% Els colors són els de la pàgina inicial (LEVELS a assets/js/materials.js);
% el logo ve de _eines/logo/pdf/ (python3 _eines/logo/genera.py).
%
% Ús, al final del preàmbul de la unitat:
%   \input{../comu/fentmates.tex}
%   \fmcurs{1eso}
% i, al començament del document:
%   \capcalera{Matemàtiques}{1r ESO}{Fitxa d'exercicis · Unitat 2}{Divisibilitat}
% Opcional abans de \input: \def\fmneed{14} (espai mínim abans de cada secció).
% =====================================================================
\makeatletter
\def\fm@logodir{../../logo/pdf}
\providecommand{\fmneed}{10}
\def\fm@cursnom{}

% clau del curs, color i color del text damunt del color
\newcommand{\fm@curs}[3]{\@namedef{fm@color@#1}{#2}\@namedef{fm@text@#1}{#3}}
\fm@curs{1eso}{E6B84A}{navy}
\fm@curs{2eso}{3B82F6}{white}
\fm@curs{3eso}{FF7A50}{white}
\fm@curs{4eso-a}{A78BFA}{white}
\fm@curs{4eso-b}{6D28D9}{white}
\fm@curs{1batx-mi}{9B1C1C}{white}
\fm@curs{1batx-ccss}{EF6B6B}{white}
\fm@curs{2batx-mii}{1E6B47}{white}
\fm@curs{2batx-ccss}{2BB3A3}{white}

\newcommand{\fmcurs}[1]{%
  \@ifundefined{fm@color@#1}{\PackageError{fentmates}{Curs desconegut: #1}{}}{}%
  \edef\fm@tmp{\noexpand\definecolor{curs}{HTML}{\@nameuse{fm@color@#1}}%
               \noexpand\colorlet{curstext}{\@nameuse{fm@text@#1}}}\fm@tmp
  \colorlet{cursfosc}{curs!72!black}%
  \def\fm@logo{\fm@logodir/vertical-#1.pdf}%
  \def\fm@marca{\fm@logodir/marca-#1.pdf}}

% ---------- capçalera de la primera pàgina ----------
% #1 matèria (petit), #2 curs, #3 tipus de document i unitat, #4 títol
\newcommand{\capcalera}[4]{%
  \gdef\fm@cursnom{#2}%
  \par\noindent
  \begin{tcolorbox}[colback=curs,colframe=curs,boxrule=0pt,arc=2.2mm,width=3.5cm,height=2.15cm,
      valign=center,left=3mm,right=2mm,top=1mm,bottom=1mm,nobeforeafter,box align=center]
    \color{curstext}{\footnotesize #1}\par\vspace{1pt}{\fontsize{19}{21}\selectfont\bfseries #2\par}%
  \end{tcolorbox}\hspace{5mm}%
  \begin{minipage}[c]{\dimexpr\linewidth-3.5cm-5mm-2.9cm\relax}
    {\color{cursfosc}\small\bfseries\MakeUppercase{#3}}\par\vspace{2pt}
    {\color{navy}\fontsize{20}{23}\selectfont\bfseries #4\par}
  \end{minipage}\hfill
  \begin{minipage}[c]{2.6cm}\raggedleft\includegraphics[width=2.6cm]{\fm@logo}\end{minipage}\par
  \vspace{2pt}}

% ---------- barres de secció del color del curs ----------
\renewcommand{\seccio}[2]{\par\needspace{\fmneed\baselineskip}\bigskip
  \begin{tcolorbox}[colback=curs,colframe=curs,coltext=curstext,arc=2mm,boxrule=0pt,left=3mm,right=3mm,top=1.2mm,bottom=1.2mm]
  {\large\bfseries #1\quad #2}
  \end{tcolorbox}\vspace{-2pt}}

% ---------- estrelles de dificultat del color del curs (a Batxillerat no se'n posen) ----------
\@ifundefined{estrelles}{}{%
  \renewcommand{\estrelles}[1]{{\color{curs}\ifcase#1\or\faStar\faStar[regular]\faStar[regular]\or\faStar\faStar\faStar[regular]\or\faStar\faStar\faStar\fi}}}

% ---------- capçalera i peu de les pàgines ----------
% primera pàgina: logos institucionals al peu
\fancypagestyle{primera}{\fancyhf{}\renewcommand{\headrulewidth}{0pt}
  \fancyfoot[L]{\raisebox{-0.32\height}{\includegraphics[height=0.78cm]{logo-gva.png}}\hspace{5mm}%
    \raisebox{-0.32\height}{\includegraphics[height=0.78cm]{logo-vermellar.png}}}
  \fancyfoot[R]{\scriptsize\color{muted}\href{https://creativecommons.org/licenses/by-nc-sa/4.0/deed.ca}{\ccbyncsa}\ \ CC BY-NC-SA 4.0 · fentmates.org\qquad\small\thepage}}
% la resta: el símbol a l'esquerra i el curs a la dreta, amb el títol al centre com abans
\fancyhead[L]{\raisebox{-0.3\height}{\includegraphics[height=0.62cm]{\fm@marca}}}
\fancyhead[R]{\small\bfseries\color{cursfosc}\fm@cursnom}
\renewcommand{\headrule}{\color{curs!45}\hrule height \headrulewidth}
\makeatother

%   \fmcurs{1eso}
% i, al començament del document:
%   \capcalera{Matemàtiques}{1r ESO}{Fitxa d'exercicis · Unitat 2}{Divisibilitat}
% Opcional abans de \input: \def\fmneed{14} (espai mínim abans de cada secció).
% =====================================================================
\makeatletter
\def\fm@logodir{../../logo/pdf}
\providecommand{\fmneed}{10}
\def\fm@cursnom{}

% clau del curs, color i color del text damunt del color
\newcommand{\fm@curs}[3]{\@namedef{fm@color@#1}{#2}\@namedef{fm@text@#1}{#3}}
\fm@curs{1eso}{E6B84A}{navy}
\fm@curs{2eso}{3B82F6}{white}
\fm@curs{3eso}{FF7A50}{white}
\fm@curs{4eso-a}{A78BFA}{white}
\fm@curs{4eso-b}{6D28D9}{white}
\fm@curs{1batx-mi}{9B1C1C}{white}
\fm@curs{1batx-ccss}{EF6B6B}{white}
\fm@curs{2batx-mii}{1E6B47}{white}
\fm@curs{2batx-ccss}{2BB3A3}{white}

\newcommand{\fmcurs}[1]{%
  \@ifundefined{fm@color@#1}{\PackageError{fentmates}{Curs desconegut: #1}{}}{}%
  \edef\fm@tmp{\noexpand\definecolor{curs}{HTML}{\@nameuse{fm@color@#1}}%
               \noexpand\colorlet{curstext}{\@nameuse{fm@text@#1}}}\fm@tmp
  \colorlet{cursfosc}{curs!72!black}%
  \def\fm@logo{\fm@logodir/vertical-#1.pdf}%
  \def\fm@marca{\fm@logodir/marca-#1.pdf}}

% ---------- capçalera de la primera pàgina ----------
% #1 matèria (petit), #2 curs, #3 tipus de document i unitat, #4 títol
\newcommand{\capcalera}[4]{%
  \gdef\fm@cursnom{#2}%
  \par\noindent
  \begin{tcolorbox}[colback=curs,colframe=curs,boxrule=0pt,arc=2.2mm,width=3.5cm,height=2.15cm,
      valign=center,left=3mm,right=2mm,top=1mm,bottom=1mm,nobeforeafter,box align=center]
    \color{curstext}{\footnotesize #1}\par\vspace{1pt}{\fontsize{19}{21}\selectfont\bfseries #2\par}%
  \end{tcolorbox}\hspace{5mm}%
  \begin{minipage}[c]{\dimexpr\linewidth-3.5cm-5mm-2.9cm\relax}
    {\color{cursfosc}\small\bfseries\MakeUppercase{#3}}\par\vspace{2pt}
    {\color{navy}\fontsize{20}{23}\selectfont\bfseries #4\par}
  \end{minipage}\hfill
  \begin{minipage}[c]{2.6cm}\raggedleft\includegraphics[width=2.6cm]{\fm@logo}\end{minipage}\par
  \vspace{2pt}}

% ---------- barres de secció del color del curs ----------
\renewcommand{\seccio}[2]{\par\needspace{\fmneed\baselineskip}\bigskip
  \begin{tcolorbox}[colback=curs,colframe=curs,coltext=curstext,arc=2mm,boxrule=0pt,left=3mm,right=3mm,top=1.2mm,bottom=1.2mm]
  {\large\bfseries #1\quad #2}
  \end{tcolorbox}\vspace{-2pt}}

% ---------- estrelles de dificultat del color del curs (a Batxillerat no se'n posen) ----------
\@ifundefined{estrelles}{}{%
  \renewcommand{\estrelles}[1]{{\color{curs}\ifcase#1\or\faStar\faStar[regular]\faStar[regular]\or\faStar\faStar\faStar[regular]\or\faStar\faStar\faStar\fi}}}

% ---------- capçalera i peu de les pàgines ----------
% primera pàgina: logos institucionals al peu
\fancypagestyle{primera}{\fancyhf{}\renewcommand{\headrulewidth}{0pt}
  \fancyfoot[L]{\raisebox{-0.32\height}{\includegraphics[height=0.78cm]{logo-gva.png}}\hspace{5mm}%
    \raisebox{-0.32\height}{\includegraphics[height=0.78cm]{logo-vermellar.png}}}
  \fancyfoot[R]{\scriptsize\color{muted}\href{https://creativecommons.org/licenses/by-nc-sa/4.0/deed.ca}{\ccbyncsa}\ \ CC BY-NC-SA 4.0 · fentmates.org\qquad\small\thepage}}
% la resta: el símbol a l'esquerra i el curs a la dreta, amb el títol al centre com abans
\fancyhead[L]{\raisebox{-0.3\height}{\includegraphics[height=0.62cm]{\fm@marca}}}
\fancyhead[R]{\small\bfseries\color{cursfosc}\fm@cursnom}
\renewcommand{\headrule}{\color{curs!45}\hrule height \headrulewidth}
\makeatother

%   \fmcurs{1eso}
% i, al començament del document:
%   \capcalera{Matemàtiques}{1r ESO}{Fitxa d'exercicis · Unitat 2}{Divisibilitat}
% Opcional abans de \input: \def\fmneed{14} (espai mínim abans de cada secció).
% =====================================================================
\makeatletter
\def\fm@logodir{../../logo/pdf}
\providecommand{\fmneed}{10}
\def\fm@cursnom{}

% clau del curs, color i color del text damunt del color
\newcommand{\fm@curs}[3]{\@namedef{fm@color@#1}{#2}\@namedef{fm@text@#1}{#3}}
\fm@curs{1eso}{E6B84A}{navy}
\fm@curs{2eso}{3B82F6}{white}
\fm@curs{3eso}{FF7A50}{white}
\fm@curs{4eso-a}{A78BFA}{white}
\fm@curs{4eso-b}{6D28D9}{white}
\fm@curs{1batx-mi}{9B1C1C}{white}
\fm@curs{1batx-ccss}{EF6B6B}{white}
\fm@curs{2batx-mii}{1E6B47}{white}
\fm@curs{2batx-ccss}{2BB3A3}{white}

\newcommand{\fmcurs}[1]{%
  \@ifundefined{fm@color@#1}{\PackageError{fentmates}{Curs desconegut: #1}{}}{}%
  \edef\fm@tmp{\noexpand\definecolor{curs}{HTML}{\@nameuse{fm@color@#1}}%
               \noexpand\colorlet{curstext}{\@nameuse{fm@text@#1}}}\fm@tmp
  \colorlet{cursfosc}{curs!72!black}%
  \def\fm@logo{\fm@logodir/vertical-#1.pdf}%
  \def\fm@marca{\fm@logodir/marca-#1.pdf}}

% ---------- capçalera de la primera pàgina ----------
% #1 matèria (petit), #2 curs, #3 tipus de document i unitat, #4 títol
\newcommand{\capcalera}[4]{%
  \gdef\fm@cursnom{#2}%
  \par\noindent
  \begin{tcolorbox}[colback=curs,colframe=curs,boxrule=0pt,arc=2.2mm,width=3.5cm,height=2.15cm,
      valign=center,left=3mm,right=2mm,top=1mm,bottom=1mm,nobeforeafter,box align=center]
    \color{curstext}{\footnotesize #1}\par\vspace{1pt}{\fontsize{19}{21}\selectfont\bfseries #2\par}%
  \end{tcolorbox}\hspace{5mm}%
  \begin{minipage}[c]{\dimexpr\linewidth-3.5cm-5mm-2.9cm\relax}
    {\color{cursfosc}\small\bfseries\MakeUppercase{#3}}\par\vspace{2pt}
    {\color{navy}\fontsize{20}{23}\selectfont\bfseries #4\par}
  \end{minipage}\hfill
  \begin{minipage}[c]{2.6cm}\raggedleft\includegraphics[width=2.6cm]{\fm@logo}\end{minipage}\par
  \vspace{2pt}}

% ---------- barres de secció del color del curs ----------
\renewcommand{\seccio}[2]{\par\needspace{\fmneed\baselineskip}\bigskip
  \begin{tcolorbox}[colback=curs,colframe=curs,coltext=curstext,arc=2mm,boxrule=0pt,left=3mm,right=3mm,top=1.2mm,bottom=1.2mm]
  {\large\bfseries #1\quad #2}
  \end{tcolorbox}\vspace{-2pt}}

% ---------- estrelles de dificultat del color del curs (a Batxillerat no se'n posen) ----------
\@ifundefined{estrelles}{}{%
  \renewcommand{\estrelles}[1]{{\color{curs}\ifcase#1\or\faStar\faStar[regular]\faStar[regular]\or\faStar\faStar\faStar[regular]\or\faStar\faStar\faStar\fi}}}

% ---------- capçalera i peu de les pàgines ----------
% primera pàgina: logos institucionals al peu
\fancypagestyle{primera}{\fancyhf{}\renewcommand{\headrulewidth}{0pt}
  \fancyfoot[L]{\raisebox{-0.32\height}{\includegraphics[height=0.78cm]{logo-gva.png}}\hspace{5mm}%
    \raisebox{-0.32\height}{\includegraphics[height=0.78cm]{logo-vermellar.png}}}
  \fancyfoot[R]{\scriptsize\color{muted}\href{https://creativecommons.org/licenses/by-nc-sa/4.0/deed.ca}{\ccbyncsa}\ \ CC BY-NC-SA 4.0 · fentmates.org\qquad\small\thepage}}
% la resta: el símbol a l'esquerra i el curs a la dreta, amb el títol al centre com abans
\fancyhead[L]{\raisebox{-0.3\height}{\includegraphics[height=0.62cm]{\fm@marca}}}
\fancyhead[R]{\small\bfseries\color{cursfosc}\fm@cursnom}
\renewcommand{\headrule}{\color{curs!45}\hrule height \headrulewidth}
\makeatother

\fmcurs{2eso}

\begin{document}
\thispagestyle{primera}

\capcalera{Matemàtiques}{2n ESO}{Fitxa d'exercicis · Unitat 6}{Proporcionalitat i percentatges}
\vspace{10pt}
Nom i cognoms: \rule{8.4cm}{0.4pt}\hfill Grup: \rule{1.6cm}{0.4pt}\hfill Data: \rule{2.2cm}{0.4pt}
\vspace{8pt}

\begin{tcolorbox}[colback=soft,colframe=soft,arc=2mm,left=3mm,right=3mm,top=1.5mm,bottom=1.5mm]
\small Exercicis graduats per nivells de dificultat:\quad \estrelles{1}\ bàsic \qquad \estrelles{2}\ mitjà \qquad \estrelles{3}\ avançat.\\[2pt] Resol-los al quadern mostrant tots els passos. Abans de calcular, decideix si la relació és directa o inversa, i comprova al final que el resultat té sentit. Les seccions segueixen l'ordre dels apunts de la unitat (6.1--6.8).
\end{tcolorbox}

% =====================================================
\seccio{1}{Raons i proporcions}

\ex{1} En un institut hi ha 360 alumnes i 24 professors. Calcula la raó entre alumnes i professors, simplifica-la i explica què significa.\recorda{La raó entre $a$ i $b$ és el quocient $\frac{a}{b}$. Una proporció és una igualtat de dues raons: $\frac{a}{b}=\frac{c}{d}\iff a\cdot d=b\cdot c$.}

\ex{1} Comprova si formen proporció:\begin{tasks}(4)
  \task $\frac{3}{4}$ i $\frac{15}{20}$
  \task $\frac{5}{6}$ i $\frac{12}{15}$
  \task $\frac{2{,}5}{2}$ i $\frac{10}{8}$
  \task $\frac{7}{9}$ i $\frac{21}{28}$
\end{tasks}

\ex{2} Calcula el terme desconegut:\begin{tasks}(4)
  \task $\frac{x}{12}=\frac{15}{20}$
  \task $\frac{8}{x}=\frac{6}{9}$
  \task $\frac{4{,}5}{6}=\frac{x}{8}$
  \task $\frac{14}{35}=\frac{6}{x}$
\end{tasks}

\ex{2} En un concurs de dibuix es presenten 84 treballs, dels quals 35 són de xics. Quina és la raó entre treballs de xics i de xiques? I entre els de xiques i el total?

\ex{3} Un batut es fa mesclant llet i maduixes en la raó $\frac{5}{2}$. Si hi posem 750 mL de llet, quantes maduixes (en mL) hi hem de posar? Quina quantitat de batut tindrem?

\ex{3} Les edats de dues germanes estan en la raó $\frac{3}{4}$ i entre les dues sumen 28 anys. Quants anys té cadascuna?


% =====================================================
\seccio{2}{Magnituds directament proporcionals}

\ex{1} Indica quines parelles de magnituds són directament proporcionals:\begin{tasks}(2)
  \task El nombre de kg de pomes i el preu.
  \task L'edat d'una persona i la seua altura.
  \task El costat d'un quadrat i el seu perímetre.
  \task Els minuts que parlem per un mòbil amb quota fixa mensual i el preu de la factura.
\end{tasks}\recorda{Són directament proporcionals si el doble d'una correspon al doble de l'altra, la meitat a la meitat… El quocient entre les quantitats corresponents és constant.}

\ex{1} Completa la taula de proporcionalitat directa i calcula'n la constant:\\[4pt]\begin{tabular}{l|ccccc}Hores & 2 & 5 & 8 &  & 20\\ \hline Km & 7 &  &  & 42 & \end{tabular}

\ex{1} Si 4 entrades de cine costen 30 €, quant costen 7 entrades? Resol-ho per reducció a la unitat.

\ex{2} Un cotxe gasta 6,4 L de gasoil cada 100 km. Resol per regla de tres:\begin{tasks}(2)
  \task Quants litres gasta en 450 km?
  \task Quants km pot fer amb 40 L?
\end{tasks}

\ex{2} Una màquina embotella 1500 botelles en 2 hores i mitja. Quantes en farà en una jornada de 8 hores? Quant tardarà a embotellar-ne 2100?

\ex{2} Per a fer una paella per a 4 persones fan falta 320 g d'arròs i 1,2 L de brou. Quant en farà falta per a 10 persones?

\ex{3} En un plànol a escala 1:250, una habitació fa 1,6 cm per 1,4 cm. Quines dimensions reals té? Quina superfície té, en m²?

\ex{3} Una empresa de lloguer de bicicletes cobra 4 € per hora. Una altra cobra 5 € fixos més 3 € per hora. a) Fes una taula per a 1, 2, 3 i 6 hores per a cada empresa. b) En quin cas el preu és directament proporcional al temps? c) A partir de quantes hores és més barata la segona?


% =====================================================
\seccio{3}{Magnituds inversament proporcionals}

\ex{1} Indica si són directament proporcionals, inversament proporcionals o cap de les dues:\begin{tasks}(2)
  \task El nombre de pintors i els dies que tarden a pintar una casa.
  \task Els litres d'aigua d'un dipòsit i el que pesa.
  \task La velocitat d'un cotxe i el temps que tarda en un trajecte.
  \task El nombre de persones d'una família i la seua despesa en menjar.
\end{tasks}\recorda{Són inversament proporcionals si el doble d'una correspon a la meitat de l'altra. El producte de les quantitats corresponents és constant.}

\ex{1} Completa la taula de proporcionalitat inversa i calcula'n la constant:\\[4pt]\begin{tabular}{l|ccccc}Aixetes & 2 & 3 & 4 &  & 12\\ \hline Hores & 18 &  &  & 6 & \end{tabular}

\ex{1} Tres màquines iguals fan una faena en 12 hores. Quant tardarien 4 màquines? Resol-ho per reducció a la unitat.

\ex{2} Un tren a 120 km/h tarda 3 h 30 min a fer un trajecte. Quant tardaria a 140 km/h?

\ex{2} Un ramader té pinso per a alimentar 40 vaques durant 30 dies. Si ven 10 vaques, quants dies li durarà el pinso? I si en compra 8 més (en lloc de vendre'n)?

\ex{2} Per a omplir un dipòsit fan falta 45 cubs de 12 L. Quants cubs de 18 L farien falta? De quina capacitat haurien de ser els cubs per a omplir-lo amb 60 cubs?

\ex{3} Unes obres s'han de fer en 18 dies i hi treballen 10 obrers. Al cap de 6 dies, se'n van 2 obrers. Quants dies tardaran els que queden a acabar l'obra?

\ex{3} Decideix si cada problema és de proporcionalitat directa o inversa i resol-lo:\begin{tasks}(1)
  \task Amb rajoles de 30 cm de costat necessitem 24 files per a cobrir una paret. Quantes files farien falta amb rajoles de 20 cm?
  \task En 5 caixes caben 120 llibres. Quantes caixes iguals fan falta per a 312 llibres?
  \task Un grup d'amics lloga una casa i, si són 6, cadascun paga 85 €. Quant pagaria cadascun si foren 10?
\end{tasks}


% =====================================================
\seccio{4}{Proporcionalitat composta}

\ex{1} Per a alimentar 6 gossos durant 10 dies calen 24 kg de pinso. Quants quilos calen per a 9 gossos durant 15 dies?\recorda{Compara cada magnitud amb la incògnita, suposant que les altres no canvien. Si és directa, multiplica per $\frac{\text{nou}}{\text{antic}}$; si és inversa, per $\frac{\text{antic}}{\text{nou}}$.}

\ex{1} 4 aixetes iguals, obertes 5 hores, aboquen 6000 L. Quants litres abocaran 6 aixetes obertes 3 hores?

\ex{2} 12 obrers, treballant 8 hores diàries, fan una obra en 20 dies. Quants dies tardaran 15 obrers treballant 8 hores diàries? I 16 obrers treballant 6 hores diàries?

\ex{2} Una família de 4 persones gasta 540 € en menjar en 30 dies. Quant gastarà una família de 6 persones en 20 dies?

\ex{2} Per a transportar 360 tones de terra, 5 camions fan 8 viatges cadascun. Quants viatges haurà de fer cada camió si en són 6 i han de transportar 540 tones?

\ex{3} En una fàbrica, 8 màquines treballant 9 hores diàries fan 3600 peces en 5 dies. Quants dies necessitaran 6 màquines treballant 10 hores diàries per a fer 4800 peces?

\ex{3} Per a llaurar un camp de 18 ha, 3 tractors treballen 6 dies. Quants tractors calen per a llaurar un camp de 30 ha en 5 dies?


% =====================================================
\seccio{5}{Repartiments proporcionals}

\ex{1} Reparteix en parts directament proporcionals:\begin{tasks}(2)
  \task 840 a 3 i 4
  \task 1200 a 2, 3 i 5
\end{tasks}\recorda{Directe: $k=\dfrac{N}{a+b+c}$ i cada part és $a\cdot k$, $b\cdot k$… Invers: es reparteix directament als inversos $\frac{1}{a},\frac{1}{b}\ldots$ (amb el m.c.m., als numeradors).}

\ex{1} Dos amics compren un dècim de loteria de 45 €: un hi posa 20 € i l'altre 25 €. Els toca un premi de 135 €. Com se l'han de repartir?

\ex{2} Tres socis han invertit 6000 €, 9000 € i 15 000 € en una empresa i han obtingut uns beneficis de 7500 €. Quant rep cadascun?

\ex{2} Reparteix en parts inversament proporcionals:\begin{tasks}(2)
  \task 500 a 2 i 3
  \task 1300 a 2, 3 i 4
\end{tasks}

\ex{2} Una àvia reparteix 1170 € entre els seus tres néts de manera directament proporcional a les seues edats: 6, 9 i 11 anys. Quant rep cadascun? I si el repartiment fóra inversament proporcional a 2, 3 i 6?

\ex{3} Un premi de 2800 € es reparteix entre tres atletes de manera inversament proporcional al lloc en què han arribat (1r, 2n i 3r). Quant rep cadascun?

\ex{3} Tres germans es reparteixen una quantitat de manera directament proporcional a 3, 5 i 7. Al segon li toquen 60 € més que al primer. Quina quantitat s'han repartit? Quant rep cadascun?


% =====================================================
\seccio{6}{Percentatges: concepte i càlcul}

\ex{1} Completa la taula:\\[4pt]\begin{tabular}{l|cccccc}Percentatge & 30\,\% & & & 2\,\% & & 150\,\% \\ \hline Fracció irreductible & & $\frac{3}{4}$ & & & $\frac{1}{8}$ & \\ Decimal & & & 0,05 & & & \end{tabular}

\ex{1} Calcula mentalment:\begin{tasks}(4)
  \task 10\,\% de 480
  \task 50\,\% de 76
  \task 25\,\% de 360
  \task 5\,\% de 240
  \task 20\,\% de 45
  \task 1\,\% de 3500
  \task 15\,\% de 60
  \task 75\,\% de 80
\end{tasks}\recorda{El 10 \% és dividir entre 10; el 5 \%, la meitat del 10 \%; el 50 \%, la meitat; el 25 \%, la quarta part; el 20 \%, la cinquena part.}

\ex{1} Calcula:\begin{tasks}(3)
  \task 12\,\% de 350
  \task 35\,\% de 84
  \task 4\,\% de 1250
  \task 7{,}5\,\% de 240
  \task 0{,}5\,\% de 6000
  \task 120\,\% de 65
\end{tasks}

\ex{2} En un institut de 850 alumnes, el 46 \% són xics, el 18 \% va a l'escola de música i el 6 \% és d'altres països. Quants alumnes hi ha en cada cas?

\ex{2} El cos humà és aproximadament un 65 \% aigua. Quants quilos d'aigua té una persona de 58 kg? I una de 72 kg?

\ex{3} D'un dipòsit de 2000 L, el dilluns se'n gasta el 30 \% i el dimarts, el 25 \% del que quedava. Quants litres queden? Quin percentatge del dipòsit s'ha gastat?


% =====================================================
\seccio{7}{Quin percentatge representa? Càlcul del total}

\ex{1} Quin percentatge representa…\begin{tasks}(2)
  \task 12 de 50?
  \task 27 de 36?
  \task 45 de 300?
  \task 90 de 72?
\end{tasks}\recorda{$p\,\%=\frac{\text{part}}{\text{total}}\cdot100$. \quad Total $=$ part $:\ \frac{p}{100}$.}

\ex{1} Calcula la quantitat total si:\begin{tasks}(2)
  \task el 20 \% és 36
  \task el 15 \% és 63
  \task el 75 \% és 240
  \task el 4 \% és 18
\end{tasks}

\ex{2} En una classe de 28 alumnes n'han aprovat 21. Quin percentatge ha aprovat? I quin ha suspés?

\ex{2} En una votació, una proposta ha obtingut 182 vots, que són el 35 \% dels vots emesos. Quantes persones han votat?

\ex{2} Una jugadora de bàsquet ha fet 15 tirs de 3 punts i n'ha encertat 6. Un company n'ha fet 20 i n'ha encertat 7. Qui té millor percentatge d'encert?

\ex{3} D'una pizza n'hem menjat el 62,5 \% i en queden 3 porcions iguals. En quantes porcions estava tallada? Quin percentatge representa cada porció?

\ex{3} Una botiga ha venut aquest mes 252 bicicletes, que són el 120 \% de les que va vendre el mes passat. Quantes en va vendre el mes passat? Quantes més n'ha venut aquest mes?


% =====================================================
\seccio{8}{Augments i disminucions percentuals}

\ex{1} Escriu l'índex de variació de cada cas:\begin{tasks}(3)
  \task augment del 25 \%
  \task disminució del 10 \%
  \task augment del 4 \%
  \task rebaixa del 35 \%
  \task IVA del 21 \%
  \task augment del 100 \%
\end{tasks}\recorda{Augment del $p\,\%$: índex $1+\frac{p}{100}$. Disminució del $p\,\%$: índex $1-\frac{p}{100}$. \ Final $=$ inicial $\cdot$ índex; \ inicial $=$ final $:$ índex.}

\ex{1} Calcula el preu final:\begin{tasks}(2)
  \task Una jaqueta de 90 € rebaixada un 30 \%.
  \task Un llibre de 22 € amb un 4 \% d'IVA.
  \task Un sou de 1450 € que puja un 3 \%.
  \task Una tauleta de 340 € amb un 15 \% de descompte.
\end{tasks}

\ex{2} Un cotxe valia 18 500 € i ha baixat un 12 \%. Quant val ara? Quants euros ha baixat?

\ex{2} Calcula la quantitat inicial:\begin{tasks}(1)
  \task Després d'una rebaixa del 20 \%, unes botes costen 68 €.
  \task Amb l'IVA del 21 \%, una factura és de 157,30 €.
  \task Després d'augmentar un 6 \%, un poble té 3180 habitants.
\end{tasks}

\ex{2} Calcula el percentatge de variació:\begin{tasks}(1)
  \task Una samarreta passa de 25 € a 20 €.
  \task Una ciutat passa de 40 000 a 43 000 habitants.
  \task El preu de l'oli passa de 6,40 € a 8 € el litre.
\end{tasks}

\ex{3} El preu d'una consola de 400 € puja un 10 \% i, després, baixa un 10 \%. Quant val al final? Quina és la variació percentual total?

\ex{3} En unes rebaixes fan un 40 \% de descompte i, a l'hora de pagar, un 10 \% més sobre el preu rebaixat. a) És el mateix que un 50 \% de descompte? b) Si hem pagat 37,80 €, quant costava abans?

\ex{3} Un comerciant compra una bicicleta per 250 €, la posa a la venda guanyant-hi un 40 \% i, com que no la ven, després la rebaixa un 20 \%. a) A quin preu la ven finalment? b) Quin percentatge guanya respecte del preu de compra?

\vfill
\begin{center}\color{muted}\small\faCheckCircle[regular]\ \ Ànim! Comprova els teus resultats amb les eines interactives dels apunts de la unitat.\end{center}

\end{document}
